% Generated draft: requires mathematical and source review.
\documentclass[11pt]{article}
\usepackage[margin=25mm]{geometry}
\usepackage{fontspec}
\setmainfont{Latin Modern Roman}
\usepackage{amsmath,amssymb,mathtools,unicode-math}
\setmathfont{Latin Modern Math}
\usepackage{enumitem,needspace,xurl,hyperref}
\setlength{\parindent}{0pt}
\setlength{\parskip}{5pt}
\setlength{\emergencystretch}{3em}
\newcommand{\N}{\mathbb{N}}
\newcommand{\Z}{\mathbb{Z}}
\newcommand{\Q}{\mathbb{Q}}
\newcommand{\R}{\mathbb{R}}
\newcommand{\C}{\mathbb{C}}
\newcommand{\F}{\mathbb{F}}
\newcommand{\A}{\mathbb{A}}
\newcommand{\PP}{\mathbb{P}}
\newcommand{\E}{\mathbb{E}}
\newcommand{\eps}{\varepsilon}
\newcommand{\dd}{\,\mathrm{d}}
\begin{document}
\section*{UnsolvedMath ID: 30000829}
\textbf{Source problem code:} catalogue entry
\subsection{Definitions and mathematical statement}
Let $(M,g)$ be the source's generic smooth Riemannian two-sphere with a metric of revolution. Poles are the fixed points of the rotational circle action. In the frequency-parameter convention, choose an orthonormal Laplace eigenbasis with
\[
-\Delta_g\phi_j=\lambda_j^2\phi_j,\qquad\lambda_j\geq0,
\]
and define
\[
N(\lambda)=\#\{j:\lambda_j\leq\lambda\},\qquad
N_{x,y}(\lambda)=\sum_{\lambda_j\leq\lambda}\phi_j(x)\overline{\phi_j(y)}.
\]
The associated global and diagonal remainders are
\[
R(\lambda)=N(\lambda)-\frac{\operatorname{Area}(M)}{4\pi}\lambda^2,
\qquad
R_x(\lambda)=N_{x,x}(\lambda)-\frac{1}{4\pi}\lambda^2.
\]
The source asks, for every $\varepsilon>0$, whether as $\lambda\to\infty$
\[
R(\lambda)=O(\lambda^{1/2+\varepsilon}),
\qquad R_x(\lambda)=O(\lambda^{1/2+\varepsilon})
\quad(x\text{ not a pole}),
\]
and
\[
N_{x,y}(\lambda)=O(\lambda^{1/2+\varepsilon})
\quad(x\ne y,\ \text{at least one of }x,y\text{ not a pole}).
\]
For fixed points, the constants may depend on the metric, $\varepsilon$ and those points; uniformity near the diagonal or poles would be a distinct stronger assertion.

The extract does not define ``generic'', specify frequency versus eigenvalue cutoff, or state uniformity of constants. The definitions above explicitly identify the standard frequency convention compatible with the displayed exponents, rather than silently identifying it with an eigenvalue cutoff. The literature assessment mentions a flat-torus specialization and the Gauss circle problem; a flat torus is not a two-sphere, so that statement is retained as a supplied comparison requiring context, not as an additional manifold admitted by the written hypothesis.
\subsection{Short English statement}
Ask for almost-square-root bounds on global and nonpolar local spectral remainders, and on off-diagonal spectral sums when at least one point is nonpolar, for generic surfaces of revolution on the sphere. Genericity, cutoff and uniformity conventions need clarification.
\subsection{Sources}
\begin{itemize}
\item Supplied report: \emph{Analysis and Geometric Singularities} (2008), Oberwolfach Reports, DOI 10.4171/owr/2007/41, pp. 2361--2442; proposed year 2007. \url{https://doi.org/10.4171/owr/2007/41}.
\item Supplied reference labelled ``Hardy--Gauss circle survey'': \url{https://arxiv.org/abs/2305.03549}.
\item Local UnsolvedMath registry, permanent ID 30000829, catalogue entry; created 25 August 2026, updated 16 September 2026. Its corrected\_verified note says the complete question and asymptotic notation were retained. The supplied 21 August 2026 literature triage labels the generic-revolution square-root remainder question open and mentions the unresolved Gauss-circle error term through a flat-torus comparison. That comparison's relation to the stated sphere hypothesis is not supplied. No independent current-status assessment is asserted here.
\end{itemize}
\end{document}
